\chapter{Incremento I:}

\section{Requisitos:}
\subsection{Requisitos funcionales:}


\begin{table}[htpb]
\centering
\caption{Requisitos Funcionales y su Incremento Asociado}
\label{tab:rf_simple}
\begin{tabularx}{\textwidth}{lXl}
\toprule
\textbf{ID} & \textbf{Descripción} & \textbf{Incremento} \\ \midrule
RF1 & Gestión integral de investigadores (CRUD). & Incremento 1 \\
RF2 & Control de periodos de estancia y fechas de finalización. & Incremento 1 \\
RF3 & Autenticación y control de acceso basado en roles (RBAC). & Incremento 1 \\
RF4 & Centralización de datos en la Fuente Única de Verdad (SSOT). & Incremento 1 \\ \midrule
RF5 & Aprovisionamiento automático de cuentas en servidores en los microservicios de CodeServer. & Incremento 2 \\
RF6 & Revocación automática de los accesos nativos. & Incremento 2 \\
RF7 & Control de acceso mediante códigos de invitación de las puertas de los laboratorios con la aplicación de NUKI & Incremento 2 \\
RF8 & Gestión de privilegios de superusuario (sudo) solo para el administrador en servidores. & Incremento 2 \\ \midrule
RF9 & Gestión de miembros en repositorios de GitLab vía API. & Incremento 3 \\
RF10 & Integración y gestión de canales en Microsoft Teams. & Incremento 3 \\ \midrule
RF11 & Registro de auditoría bajo el estándar \textbf{NIST SP 800-92r1}. & Incremento 4 \\ \bottomrule
\end{tabularx}
\end{table}

\subsection{Requisitos no funcionales:}
